% ============================================================
%  CHAPITRE I : Contexte général et problématique
% ============================================================
\chapter{Contexte général et problématique}
\label{chap:contexte}

\section*{Introduction}

Ce chapitre situe le cadre du projet. Il présente l'organisme d'accueil et son
organisation logistique, décrit la situation actuelle de la gestion des stocks,
formule la problématique retenue et précise le cahier des charges ainsi que la
planification adoptée.

% -------------------------------------------------------
\section{Présentation de l'organisme d'accueil}
\label{sec:organisme}

\subsection{Le groupe OCP}

\hspace{1cm}[Présenter le groupe \gls{ocp} : historique, métiers (extraction,
valorisation, commercialisation des phosphates et dérivés), implantations
géographiques, chiffres clés et position sur le marché mondial. Citer les
sources utilisées, par exemple le rapport annuel du groupe.]

\subsection{Le site d'accueil et l'entité de rattachement}

\hspace{1cm}[Décrire le site sur lequel s'est déroulé le stage ainsi que
l'entité d'accueil : sa mission, son périmètre et sa place dans l'organisation
du groupe.]

\subsection{Structure organisationnelle}

La figure~\ref{fig:organigramme} présente la position de l'entité d'accueil au
sein de l'organisation.

\begin{figure}[H]
  \centering
  \begin{tikzpicture}[
    node distance=1.5cm,
    box/.style={
      rectangle, draw=iacsTeal, fill=iacsLight!40,
      rounded corners=4pt, minimum width=3.5cm,
      minimum height=0.9cm, align=center,
      font=\small\sffamily
    },
    arrow/.style={-Stealth, thick, color=iacsTeal}
  ]
    \node[box, fill=iacsTeal!20, font=\small\sffamily\bfseries] (dg) {Direction du site};
    \node[box, below left=1.2cm and 1.2cm of dg] (achat) {Achats};
    \node[box, below=1.2cm of dg] (log) {Supply Chain\\/ Logistique};
    \node[box, below right=1.2cm and 1.2cm of dg] (maint) {Maintenance};
    \node[box, below=1.2cm of log, fill=iacsGreen!20] (mag) {Gestion des stocks\\(magasins)};

    \draw[arrow] (dg) -- (achat);
    \draw[arrow] (dg) -- (log);
    \draw[arrow] (dg) -- (maint);
    \draw[arrow] (log) -- (mag);
  \end{tikzpicture}
  \caption{Positionnement de l'entité d'accueil dans l'organisation}
  \label{fig:organigramme}
\end{figure}

\iacsnote{Remplacer cet organigramme par l'organigramme réel communiqué par
l'encadrant industriel.}

% -------------------------------------------------------
\section{Le processus de gestion des stocks}
\label{sec:processus}

\subsection{Description du flux actuel}

La figure~\ref{fig:flux_appro} schématise le cycle de réapprovisionnement, de
l'expression du besoin jusqu'à la mise à disposition de l'article.

\begin{figure}[H]
  \centering
  \begin{tikzpicture}[
    node distance=0.9cm and 1.2cm,
    step/.style={
      rectangle, draw=iacsTeal, fill=iacsLight!50,
      rounded corners=4pt, minimum width=2.6cm,
      minimum height=1cm, align=center,
      font=\scriptsize\sffamily
    },
    arrow/.style={-Stealth, thick, color=iacsTeal}
  ]
    \node[step] (besoin)   {Expression\\du besoin};
    \node[step, right=of besoin]  (da)   {Demande\\d'achat};
    \node[step, right=of da]      (cmd)  {Commande\\fournisseur};
    \node[step, below=1.2cm of cmd] (recep) {Réception\\\& contrôle};
    \node[step, left=of recep]    (stock) {Mise en\\stock};
    \node[step, left=of stock]    (sortie){Sortie pour\\consommation};

    \draw[arrow] (besoin) -- (da);
    \draw[arrow] (da) -- (cmd);
    \draw[arrow] (cmd) -- (recep);
    \draw[arrow] (recep) -- (stock);
    \draw[arrow] (stock) -- (sortie);
  \end{tikzpicture}
  \caption{Cycle de réapprovisionnement et de consommation}
  \label{fig:flux_appro}
\end{figure}

\subsection{Système d'information et données disponibles}

\hspace{1cm}[Préciser l'outil de gestion utilisé (\gls{erp} de type \gls{sap} ou
autre), les données extraites, leur profondeur d'historique et leur granularité.]

\subsection{Diagnostic de la situation actuelle}

\hspace{1cm}[Restituer ici le diagnostic établi à partir des entretiens et des
premières analyses de données : niveau de service constaté, valeur du stock,
taux de rotation, part des références dormantes, respect des délais
fournisseurs.]

% -------------------------------------------------------
\section{Problématique et objectifs}
\label{sec:problematique}

\subsection{Problématique}

\hspace{1cm}La gestion des stocks fait face à un arbitrage structurel. Un stock
insuffisant expose à des \textbf{ruptures}, dont le coût réel n'est pas celui de
l'article manquant mais celui de l'arrêt d'exploitation qu'il provoque. À
l'inverse, un stock excédentaire immobilise du capital, occupe de l'espace et
expose au risque d'obsolescence.

Cet arbitrage est compliqué par deux sources d'incertitude : la
\textbf{variabilité de la consommation} et celle des \textbf{délais
fournisseurs}. En l'absence de règles de réapprovisionnement dimensionnées sur
ces variabilités, les niveaux de stock résultent le plus souvent de décisions
empiriques, hétérogènes d'une référence à l'autre.

La question posée est donc la suivante : \emph{comment dimensionner les
paramètres de gestion des stocks de manière à garantir un taux de service cible
tout en minimisant la valeur du stock immobilisé ?}

\iacsnote{Étayer la problématique par des chiffres issus des données réelles
(nombre de ruptures constatées, valeur du stock dormant, écart entre délai
théorique et délai réel).}

\subsection{Objectifs du projet}

Les objectifs principaux de ce projet sont :

\begin{itemize}
  \item Constituer et fiabiliser une base de données exploitable à partir des
    historiques de consommation et de stock.
  \item Segmenter le portefeuille de références selon leur valeur
    (classification \gls{abc}) et la régularité de leur demande
    (classification \gls{xyz}).
  \item Caractériser statistiquement la demande et les délais de
    réapprovisionnement de chaque segment.
  \item Calculer les paramètres de gestion (stock de sécurité, point de
    commande, quantité de commande) adaptés à chaque classe.
  \item Quantifier les gains attendus en taux de service et en valeur de stock.
  \item Livrer un outil d'aide à la décision exploitable par les équipes.
\end{itemize}

\subsection{Cahier des charges}

Le tableau~\ref{tab:cahier_charges} résume les spécifications retenues.

\begin{table}[H]
  \centering
  \caption{Cahier des charges du projet}
  \label{tab:cahier_charges}
  \renewcommand{\arraystretch}{1.4}
  \begin{tabularx}{\textwidth}{|>{\bfseries\color{iacsPrimary}}l|X|}
    \hline
    \rowcolor{iacsTeal}
    \textbf{\textcolor{white}{Critère}} & \textbf{\textcolor{white}{Description}} \\
    \hline
    Périmètre & [Familles ou références retenues, à arrêter avec l'encadrant] \\
    \hline
    Données d'entrée & Historique des consommations et des stocks sur 2 à 3 ans, au format Excel \\
    \hline
    Fonctionnalité 1 & Classification automatique des références (\gls{abc}/\gls{xyz}) \\
    \hline
    Fonctionnalité 2 & Calcul du stock de sécurité et du point de commande par référence \\
    \hline
    Fonctionnalité 3 & Simulation de l'impact d'un taux de service cible sur la valeur du stock \\
    \hline
    Fonctionnalité 4 & Tableau de bord interactif et export des résultats sous Excel \\
    \hline
    Contrainte & Solution reposant exclusivement sur des outils libres et gratuits \\
    \hline
    Livrable & Rapport, code source versionné et tableau de bord opérationnel \\
    \hline
  \end{tabularx}
\end{table}

% -------------------------------------------------------
\section{Planification du projet}
\label{sec:planification}

La figure~\ref{fig:gantt} illustre le déroulement prévisionnel du projet.

\begin{figure}[H]
  \centering
  \begin{tikzpicture}[x=0.9cm, y=0.7cm]
    % --- Fond alterné ---
    \fill[iacsLightGray] (0,-0.5) rectangle (12,0.5);
    \fill[white]         (0,-1.5) rectangle (12,-0.5);
    \fill[iacsLightGray] (0,-2.5) rectangle (12,-1.5);
    \fill[white]         (0,-3.5) rectangle (12,-2.5);
    \fill[iacsLightGray] (0,-4.5) rectangle (12,-3.5);
    \fill[white]         (0,-5.5) rectangle (12,-4.5);

    % --- Barres du Gantt ---
    \fill[iacsTeal,          rounded corners=2pt] (0.2,0.1)   rectangle (2.2,0.4);
    \fill[iacsAccent,        rounded corners=2pt] (1.8,-0.9)  rectangle (4.5,-0.6);
    \fill[iacsGreen!80!black,rounded corners=2pt] (4.2,-1.9)  rectangle (7.0,-1.6);
    \fill[iacsOrange,        rounded corners=2pt] (6.5,-2.9)  rectangle (9.0,-2.6);
    \fill[iacsRed!80,        rounded corners=2pt] (8.5,-3.9)  rectangle (10.5,-3.6);
    \fill[iacsSecondary,     rounded corners=2pt] (2.0,-4.9)  rectangle (11.8,-4.6);

    % --- Labels ---
    \node[anchor=east, font=\small\sffamily] at (-0.2, 0) {Cadrage \& entretiens};
    \node[anchor=east, font=\small\sffamily] at (-0.2,-1) {Collecte \& nettoyage};
    \node[anchor=east, font=\small\sffamily] at (-0.2,-2) {Analyse \& segmentation};
    \node[anchor=east, font=\small\sffamily] at (-0.2,-3) {Modélisation \& calculs};
    \node[anchor=east, font=\small\sffamily] at (-0.2,-4) {Tableau de bord};
    \node[anchor=east, font=\small\sffamily] at (-0.2,-5) {Rédaction du rapport};

    % --- Axe semaines ---
    \foreach \m/\n in {0/S1, 2/S2, 4/S3, 6/S4, 8/S5, 10/S6} {
      \node[anchor=north, font=\scriptsize\sffamily\color{iacsGray}] at (\m+1, 1) {\n};
    }
    \draw[iacsGray, thin] (0,0.7) -- (12,0.7);
  \end{tikzpicture}
  \caption{Diagramme de Gantt prévisionnel du projet}
  \label{fig:gantt}
\end{figure}

\section*{Conclusion}

Ce chapitre a situé le cadre du projet, décrit le processus de gestion des
stocks en place et formulé la problématique retenue. Le chapitre suivant
présente les fondements théoriques et les méthodes mobilisées pour y répondre.
